\chapter[Lezione 5 --- 6 ottobre 2026]{Lezione 5: disuguaglianza di Bernoulli, binomio di Newton ed esercizi sugli estremi}
\label{chap:lezione-05}

\section{Disuguaglianza di Bernoulli}
\label{sec:lez05-bernoulli}

\begin{theorem}[Disuguaglianza di Bernoulli]
\label{thm:disuguaglianza-bernoulli}
Per ogni \(x\in\R\) con \(x>-1\) e per ogni \(n\in\N\) vale
\[
  (1+x)^n\ge 1+nx.
\]
\end{theorem}

\begin{proof}
Sia \(x>-1\) fissato. Si procede per induzione su \(n\),
applicando il \cref{thm:principio-induzione} con valore iniziale
\(p=0\).

\emph{Passo base.}
Per \(n=0\), i due membri coincidono:
\[
  (1+x)^0=1=1+0\cdot x.
\]

\emph{Passo induttivo.}
Sia \(n\in\N\) e si supponga che
\[
  (1+x)^n\ge 1+nx.
\]
Poiché \(x>-1\), si ha \(1+x>0\). Moltiplicando per
\(1+x\), il verso della disuguaglianza rimane invariato:
\begin{align*}
  (1+x)^{n+1}
    &\ge (1+nx)(1+x)\\
    &=1+x+nx+nx^2\\
    &=1+(n+1)x+nx^2\\
    &\ge 1+(n+1)x.
\end{align*}
L'ultima disuguaglianza segue da \(nx^2\ge0\), poiché
\(n\ge0\) e \(x^2\ge0\). Il passo induttivo è quindi verificato.

Per il principio di induzione, la disuguaglianza vale per ogni
\(n\in\N\). Poiché \(x>-1\) era arbitrario, segue la tesi.
\end{proof}

\section{Binomio di Newton}
\label{sec:lez05-binomio-newton}

I coefficienti binomiali introdotti nella
\cref{def:coefficiente-binomiale} permettono di esprimere la potenza
di una somma mediante una sommatoria finita.

\begin{theorem}[Binomio di Newton]
\label{thm:binomio-newton}
Per ogni \(a,b\in\R\) e per ogni \(n\in\N\) vale
\[
  (a+b)^n=\sum_{k=0}^{n}\binom{n}{k}a^k b^{n-k}.
\]
Si adotta la convenzione \(t^0=1\), anche per \(t=0\).
\end{theorem}

Per \(n\ge4\), lo sviluppo assume la forma
\[
  (a+b)^n
  =b^n+nab^{n-1}+\cdots+na^{n-1}b+a^n.
\]
In particolare, per \(n=2\) e \(n=3\) si ottengono le identità
\[
  (a+b)^2=b^2+2ab+a^2,
\]
\[
  (a+b)^3=b^3+3ab^2+3a^2b+a^3.
\]

Il teorema si dimostra mediante il \cref{thm:principio-induzione},
utilizzando l'identità di Pascal della \cref{prop:identita-pascal}.
La dimostrazione completa è riportata
nell'\cref{ex:tutorato01-binomio-newton} del primo tutorato.

\section{Esercizi sul fattoriale}
\label{sec:lez05-esercizi-fattoriale}

\begin{exercise}[Confronto tra fattoriale e potenza di \(2\)]
\label{ex:fattoriale-maggiore-due-n}
Dimostrare che, per ogni \(n\in\N\) con \(n\ge4\), vale
\[
  n!>2^n.
\]
\end{exercise}

\begin{proof}[Soluzione dell'esercizio]
Per \(n<4\) si ha
\[
\begin{aligned}
  0!&=1=2^0,\\
  1!&=1<2=2^1,\\
  2!&=2<4=2^2,\\
  3!&=6<8=2^3.
\end{aligned}
\]
La disuguaglianza richiesta non vale quindi per \(n<4\).
Si procede per induzione, applicando il
\cref{thm:principio-induzione} con valore iniziale \(p=4\).

\emph{Passo base.}
Per \(n=4\),
\[
  4!=24>16=2^4.
\]

\emph{Passo induttivo.}
Sia \(n\ge4\) e si supponga \(n!>2^n\).
Per la \cref{def:fattoriale}, moltiplicando per \(n+1>0\),
si ottiene
\[
  (n+1)!=(n+1)n!>(n+1)2^n.
\]
Poiché \(n+1\ge5>2\), si ha inoltre
\[
  (n+1)2^n>2\cdot2^n=2^{n+1}.
\]
Pertanto \((n+1)!>2^{n+1}\), e il passo induttivo è verificato.

Per il \cref{thm:principio-induzione}, la disuguaglianza
vale per ogni \(n\ge4\).
\end{proof}

\begin{exercise}[Confronto tra fattoriale e quadrato]
\label{ex:fattoriale-maggiore-n-quadro}
Dimostrare che, per ogni \(n\in\N\) con \(n\ge4\), vale
\[
  n!>n^2.
\]
\end{exercise}

\begin{proof}[Soluzione dell'esercizio]
Nel passo induttivo sarà necessario confrontare \(n^2\)
con \(n+1\). Per ogni \(n\ge2\) si ha
\[
  n^2-(n+1)=n(n-1)-1\ge2-1=1>0,
\]
quindi \(n^2>n+1\).

Si procede per induzione, applicando il
\cref{thm:principio-induzione} con valore iniziale \(p=4\).

\emph{Passo base.}
Per \(n=4\),
\[
  4!=24>16=4^2.
\]

\emph{Passo induttivo.}
Sia \(n\ge4\) e si supponga \(n!>n^2\).
Per la \cref{def:fattoriale}, moltiplicando per \(n+1>0\),
si ottiene
\[
  (n+1)!=(n+1)n!>(n+1)n^2.
\]
Poiché \(n\ge4\), dalla disuguaglianza appena dimostrata
segue \(n^2>n+1\). Pertanto
\[
  (n+1)!>(n+1)n^2>(n+1)^2.
\]
Il passo induttivo è verificato. Per il
\cref{thm:principio-induzione}, la disuguaglianza
vale per ogni \(n\ge4\).
\end{proof}

\section{Esercizi sugli estremi}
\label{sec:lez05-esercizi-estremi}

Negli esercizi si utilizzano le
\cref{def:estremo-superiore,def:estremo-inferiore} e le
convenzioni per gli estremi infiniti della \cref{sec:lez04-estremi},
con l'ordinamento della \cref{def:lez07-retta-estesa}.

\begin{exercise}[Estremi e inclusione]
\label{ex:estremi-inclusione}
Siano \(A,B\subseteq\R\) non vuoti, con \(A\subseteq B\).
\begin{enumerate}
  \item Stabilire se necessariamente \(\sup A\in B\).
  \item Dimostrare che \(\inf B\le\inf A\).
  \item Dimostrare che \(\sup A\le\sup B\).
\end{enumerate}
\end{exercise}

\begin{proof}[Soluzione dell'esercizio]
\emph{Primo punto.}
L'appartenenza \(\sup A\in B\) non è garantita, neppure quando
\(A\) è superiormente limitato. Si considerino, ad esempio,
\[
  A=(0,1),\qquad B=(0,1)\cup\{2\}.
\]
Si ha \(A\subseteq B\), ma \(\sup A=1\notin B\):
l'estremo superiore di \((0,1)\) è \(1\) per il
\cref{thm:caratterizzazione-sup}.
Se \(A\) non è superiormente limitato, \(\sup A=+\infty\)
non può appartenere a \(B\subseteq\R\).

\emph{Secondo punto.}
Se \(B\) non è inferiormente limitato, si ha \(\inf B=-\infty\),
quindi \(\inf B\le\inf A\) nell'ordinamento esteso.

Se \(B\) è inferiormente limitato, per la
\cref{prop:limitatezza-sottoinsiemi} lo è anche \(A\).
Per la \cref{def:estremo-inferiore}, \(\inf B\) è un minorante
di \(B\) e, poiché \(A\subseteq B\), si ha
\[
  \inf B\le a\qquad\forall a\in A.
\]
Dunque \(\inf B\in m_A\). Essendo \(\inf A=\max m_A\),
si conclude
\[
  \inf B\le\inf A.
\]

\emph{Terzo punto.}
Se \(B\) non è superiormente limitato, si ha \(\sup B=+\infty\),
quindi \(\sup A\le\sup B\) nell'ordinamento esteso.

Se \(B\) è superiormente limitato, per la
\cref{prop:limitatezza-sottoinsiemi} lo è anche \(A\).
Per la \cref{def:estremo-superiore}, \(\sup B\) è un maggiorante
di \(B\) e, poiché \(A\subseteq B\), si ha
\[
  a\le\sup B\qquad\forall a\in A.
\]
Dunque \(\sup B\in M_A\). Essendo \(\sup A=\min M_A\),
si conclude
\[
  \sup A\le\sup B.
\]
\end{proof}

\begin{exercise}[Estremi dell'insieme degli opposti]
\label{ex:estremi-insieme-opposto}
Sia \(A\subseteq\R\) non vuoto e si consideri l'insieme
\(-A\) definito nella \cref{sec:lez02-numeri-interi}.
Dimostrare che
\[
  \inf(-A)=-\sup A,
  \qquad
  \sup(-A)=-\inf A.
\]
Per gli estremi infiniti si adottano le convenzioni
\(-(+\infty)=-\infty\) e \(-(-\infty)=+\infty\).
\end{exercise}

\begin{proof}[Soluzione dell'esercizio]
Poiché \(A\) è non vuoto, lo è anche \(-A\).
Cambiando segno, il verso delle disuguaglianze si inverte:
\(a\le t\Longleftrightarrow-a\ge-t\).

\emph{Prima identità.}
Si supponga che \(A\) sia superiormente limitato e si ponga
\(s=\sup A\in\R\). Per la \cref{def:estremo-superiore},
si ha \(a\le s\) per ogni \(a\in A\), quindi
\[
  -s\le-a\qquad\forall a\in A.
\]
Pertanto \(-s\) è un minorante di \(-A\).

Sia \(\ell\in\R\) un qualsiasi minorante di \(-A\).
Allora \(\ell\le-a\) per ogni \(a\in A\), da cui
\[
  a\le-\ell\qquad\forall a\in A.
\]
Il numero \(-\ell\) è dunque un maggiorante di \(A\).
Poiché \(s\) è il minimo dei maggioranti, si ha
\(s\le-\ell\), ossia \(\ell\le-s\).
Pertanto \(-s\) è il massimo dei minoranti di \(-A\), e
per la \cref{def:estremo-inferiore} si conclude
\[
  \inf(-A)=-s=-\sup A.
\]

Se \(A\) non è superiormente limitato, per ogni \(t\in\R\)
esiste \(a\in A\) tale che \(a>-t\); altrimenti \(-t\)
sarebbe un maggiorante di \(A\). Cambiando segno, si ottiene
\[
  -a<t.
\]
L'insieme \(-A\) non è quindi inferiormente limitato.
Poiché \(\sup A=+\infty\), si conclude
\[
  \inf(-A)=-\infty=-\sup A.
\]

\emph{Seconda identità.}
Si supponga che \(A\) sia inferiormente limitato e si ponga
\(i=\inf A\in\R\). Per la \cref{def:estremo-inferiore},
si ha \(i\le a\) per ogni \(a\in A\), quindi
\[
  -a\le-i\qquad\forall a\in A.
\]
Pertanto \(-i\) è un maggiorante di \(-A\).

Sia \(u\in\R\) un qualsiasi maggiorante di \(-A\).
Allora \(-a\le u\) per ogni \(a\in A\), da cui
\[
  -u\le a\qquad\forall a\in A.
\]
Il numero \(-u\) è dunque un minorante di \(A\).
Poiché \(i\) è il massimo dei minoranti, si ha
\(-u\le i\), ossia \(-i\le u\).
Pertanto \(-i\) è il minimo dei maggioranti di \(-A\), e
per la \cref{def:estremo-superiore} si conclude
\[
  \sup(-A)=-i=-\inf A.
\]

Se \(A\) non è inferiormente limitato, per ogni \(t\in\R\)
esiste \(a\in A\) tale che \(a<-t\); altrimenti \(-t\)
sarebbe un minorante di \(A\). Cambiando segno, si ottiene
\[
  -a>t.
\]
L'insieme \(-A\) non è quindi superiormente limitato.
Poiché \(\inf A=-\infty\), si conclude
\[
  \sup(-A)=+\infty=-\inf A.
\]
\end{proof}
